% DRAFT of the new subsection of BALABAN_SUPPLEMENT.tex (after sec:far-terms). Not a deliverable and not a source.
% Sources: W-delta note 27210707d749, #W11.136--#W11.143 and Section 12 (#W12.1--#W12.257); LF note ada99ead6f26.
% Each part is written when its part of the note has been read and its citations of Balaban re-read.
% Parts written: 1 (the glued configuration at the cut), 5 (the gauge transformations). Parts to write: 2--4, 6.

\subsection{The Localisation Functions of an $\mathbf R$-Operation on Two Fine Lattices}\label{sec:localisation-functions}

The expansions of an $\mathbf R$-operation are formed with three functions which Lemma~\ref{lem:field-difference} does not cover as it stands. We recall them. Ba{\l}aban obtains the configuration $U''_k$ from the new background $U_k$ in four steps of the form $U' = (\exp i\eta\mathbf H\,U)^{u^{-1}}$ \cite[(1.51), (1.54), (1.57), (1.58), pp.~370--371]{BalabanLFR2}: the function $\mathbf H_k(-i\log V'_k\lceil_\Lambda)$ of the set $\mathbf B_k$ about $U_k$; the function $\mathbf H_{\mathbf B_1}(-i\log M^\cdot(U_{k,\Lambda})\lceil_{Z''_k})$ of the set $\mathbf B_1 = \mathbf B''_k \cup \mathbf B_h(\Omega''^\sim_{h+1})$ about $U_{k,\Lambda}$; the function $\mathbf H''_k(\mathbf G''_kJ_{1,2}, H''_{1,k}\tilde Q_\cdot(\eta\mathbf A_0))$ of the set $\mathbf B''_k$ about the configuration $U_{1,2}$, which is $U_1$ on $\Omega''^\sim_{h+1}$ and $U_2$ on its complement \cite[(1.52), (1.53)]{BalabanLFR2}; and the function $\mathbf H''_k(g_kB)$ of the same set about $U^0_k$. In the localisation expansions \cite[(1.61)--(1.64), pp.~373--374]{BalabanLFR2} the first three are taken with the parameters $s(\Delta)$ of the cubes $\Delta$ which are disjoint from $\Lambda$, and each of them is accompanied by a gauge transformation, which is ``an analytic function of'' the function ``restricted to this block'', ``given by the explicit formula (106) [12]'' \cite[p.~373]{BalabanLFR2}. The first and the fourth function are of the form \cite[(179), (180)]{BalabanVAR}, the expansion of a minimiser of a determining set about the minimiser of the same set at other data, and Lemma~\ref{lem:field-difference} compares them between the two runs. The second and the third are not. The background $U_{1,2}$ of the third is not a critical configuration of $\mathbf B''_k$, and Ba{\l}aban takes the function from \cite[Sects.~D, E]{BalabanVAR}, with the current of the background among its arguments; the set $\mathbf B_1$ of the second has, along the boundary of $\Omega''^\sim_{h+1}$, a collar whose layers descend to the fine lattice \cite[p.~256, ll.~57--62; (2.14), p.~257]{BalabanCRE}, and there the lattices of the two runs differ. This subsection compares the three functions and their gauge transformations between the two runs. Throughout, as in \cite[p.~177, ll.~33--35]{BalabanLFR1}, the domains of \cite[(1.50)--(1.55)]{BalabanLFR2} are their intersections with $Z$, so that $U_1$ lives on $\Omega''^\sim_{h+1} \cup Z^c$, $U_2$ on $(\Omega''^\sim_{h+1})^c \cap Z$, and the only interface between them is $\partial\Omega''^\sim_{h+1} \cap Z$; with the complements taken in the whole lattice the statements ``considered on $Z^c$'' and ``restricted to $Z^c$ is equal to $U_1$'' of \cite[p.~370, ll.~14--15; p.~371, ll.~26--27]{BalabanLFR2} would not hold.

\begin{lemma}[The Glued Configuration at the Cut]\label{lem:glued-configuration}
Let $\tilde\Omega \coloneqq \Omega''^\sim_{h+1}$, let the \emph{band} be the set $(\Omega''_{h+1})^c \cap \Omega''^{\sim2}_{h+1} \cap Z$, and let $\varepsilon_U$ bound, in the sense of \cite[(7)]{BalabanVAR}, the data of $U_1$ and of $U_2$ at a real point of the region of integration. Let the block size $M'$ of Lemma~\ref{lem:joins-block-size} satisfy
\begin{equation}\label{eq:band-condition}
    B_3e^{-\frac{1}{2}\delta_0M'} \leq \tfrac{1}{2}L^{-3}\;,
\end{equation}
with the constants $B_3$ and $\delta_0$ of \cite[(162), p.~303]{BalabanVAR} at this block size, and let $\gamma$ be so small that $4K_{\mathrm F}M'^2B_3\varepsilon_U \leq \frac{1}{2}L^{-3}$, where $K_{\mathrm F} \coloneqq B_0(C_4 + 4C_2)(36dL^2B_1)^2$ is the coefficient of \cite[(165), p.~304]{BalabanVAR}. Then, in either run, the band lies in the layer $h$ of $\mathbf B''_k$, and $U_{1,2}$ satisfies there the conditions \cite[(2), p.~278]{BalabanVAR} of the sequence of $\mathbf B''_k$, at the index $h$, with $L^3B_3\varepsilon_U$ in the place of $\varepsilon_0$; outside the band it satisfies them with $B_3\varepsilon_U$. On the bonds within one bond of $\partial\tilde\Omega \cap Z$ it is the trivial configuration. Consequently $U_{1,2}$ satisfies the hypothesis \cite[(14), p.~280]{BalabanVAR} of the base of \cite[Sects.~D, E]{BalabanVAR} for the set $\mathbf B''_k$, with the constant $C_1 \coloneqq L^3$, and the conditions \cite[(3.35), (3.36)]{BalabanPBG}.
\end{lemma}

\begin{proof}
Ba{\l}aban states the regularity near the cut and gives its reason in one sentence: the configuration ``is regular, and satisfies the regularity conditions typical for an $h$-order configuration on the domain $(\Omega''_{h+1})^c \cap \Omega''^{\sim2}_{h+1}$. It is so, because the arguments of the corresponding functions are equal to $1$ on this domain, and the exponential decay properties enforce the additional regularity'' \cite[p.~370, ll.~35--38]{BalabanLFR2}. We carry the reason out.

\emph{The band and the strips.} On $Z$, $\Omega''_{h+1} = \Omega^{\sim5}_{h+1} \cap Z$ \cite[(1.11), (1.12), p.~179]{BalabanLFR1}, so that the band lies in $\Omega^{\sim7}_{h+1} \setminus \Omega_{h+1}$, and since $\dist(\Omega_{h+1}, \Omega_h^c) \geq 8MR_{h+1}a_{h+1}$ (Lemma~\ref{lem:layer-widths}) it lies in $\Omega''_h \setminus \Omega''_{h+1}$, the layer $h$ of $\mathbf B''_k$. By \cite[(1.52)]{BalabanLFR2} and \cite[(2.14)]{BalabanCRE} the set $\mathbf B_1$ carries on $\tilde\Omega \setminus \tilde\Omega^{\sim-2}$ the set $\mathbf B_h(\tilde\Omega)$, a sequence $\tilde\Omega = \Omega_0 \supset \Omega_1 \supset \dots \supset \Omega_h$ of unions of cubes of the sides $M_1a_n$ with $\dist(\Omega_n, \Omega_{n-1}^c) \geq M_1a_n$; its finest layer is a strip along $\partial\tilde\Omega$ on which the variables of the determining set are the bonds themselves \cite[(3), (4), p.~278]{BalabanVAR}. The set $\mathbf B_2$ has the same structure on the other side. Both collars lie within $2M_1a_h$ of $\partial\tilde\Omega$ \cite[p.~256, l.~62]{BalabanCRE}, hence in the band, which reaches one layer of cubes of the side $MR_{h+1}a_{h+1}$ on either side. The band lies in $Z''_{h+1} \subset Z''_k$ \cite[p.~195, l.~42]{BalabanLFR1}, and by \cite[(1.53)]{BalabanLFR2} the data of $U_1$ on its part in $\tilde\Omega$ and those of $U_2$ on its part in $\tilde\Omega^c$ are equal to $1$; the other data of $U_1$, the averages $M^k(U^{(AL)}_0)$ on $\Lambda \cap \Omega''_k$ and the field $V''$ on $\Lambda^c$, lie off the band, because $\Omega''_k \subset \Omega''_{h+1}$ and $\Lambda^c \subset (Z''_k)^c$.

\emph{The cut.} A configuration whose bonds on the two finest strips are prescribed to be $1$ is $1$ there, whatever representative of the minimal orbit is taken. The boundary $\partial\tilde\Omega$ is a union of faces of cubes, so that every plaquette lies in the closure of $\tilde\Omega$ or in that of its complement, and the bonds on $\partial\tilde\Omega$ are $1$ from both sides. A plaquette variable or a current which reads bonds of both $U_1$ and $U_2$ reads bonds within one bond of $\partial\tilde\Omega$ only; there the plaquette variables of $U_{1,2}$ are $1$ and its current vanishes, and no term which crosses the cut remains.

\emph{Off the strips.} Outside the collars the layers of $\mathbf B_1$ and of $\mathbf B_2$ are those of $\mathbf B''_k$, and \cite[Theorem~1]{BalabanVAR} for $U_1$ and $U_2$, at the block size of Lemma~\ref{lem:joins-block-size}, gives \cite[(2)]{BalabanVAR} with $B_3\varepsilon_U$.

\emph{On the strips.} At a point of a layer $n < h$ of a collar, \cite[Theorem~1]{BalabanVAR} gives the conditions (2) at the index $n$ only, and at the index $h$ they are stronger by the factor $L^{2(h-n)}$ for the plaquette variables and $L^{3(h-n)}$ for the currents. The gain comes from the proof of that theorem. One step of \cite[Sect.~F, pp.~300--304]{BalabanVAR}, at a cube $\square_k$ of a layer with the enlarged cube $\tilde\square$, bounds the linear part by $\frac{1}{4}\max\{B_3p^{\mathrm{loc}}, B_3e^{-\frac{1}{2}\delta_0R_1M_1}e^{\mathrm{old}}\}$, with $p^{\mathrm{loc}}$ the bound of the plaquette variables of the data on $\square_k$ and $e^{\mathrm{old}}$ the supremum over $\tilde\square$ of the bounds (2), each at its own layer \cite[(161), (162)]{BalabanVAR}; the nonlinear part adds $K_{\mathrm F}(R_1M_1)^2(e^{\mathrm{old}})^2$ \cite[(165)]{BalabanVAR}; and \cite[(168)]{BalabanVAR} doubles the result. With $\frac{1}{2}x + y \leq \max(x, 2y)$, one step gives (2) on the inner cube with $\max\{B_3p^{\mathrm{loc}}, \theta e^{\mathrm{old}}\}$, where $\theta \coloneqq B_3e^{-\frac{1}{2}\delta_0M'} + 4K_{\mathrm F}M'^2e^{\mathrm{old}}$ and $R_1M_1 = M'$ at the joins. Under \eqref{eq:band-condition} and the smallness of $\varepsilon_U$, $\theta \leq L^{-3}$ at every step, since the bounds only decrease. Iterating, the bound at a block is at most the supremum over the chains of steps of $\theta^{\mathrm{length}}B_3p^{\mathrm{loc}}$ at the end of the chain, the term $\theta^m$ times the starting bound tending to zero. A step at a block of a layer $n$ reads the cube $\tilde\square$, which lies within $6dM'a_n$ of it \cite[p.~300, ll.~28--37]{BalabanVAR}, and the layer $n + 1$ of the collar is at least $LM_1a_{n+1}$ wide, with $M_1 = K_{M'}M' \geq (12d + 1)M'$; so $\tilde\square$ does not meet the layer $n + 2$, and one step changes the layer by at most one. On the strips and on the band the data are $1$, so that $p^{\mathrm{loc}} = 0$ there; the other data lie on blocks of $\mathbf B''_k$ off the band, which a chain from a block of the layer $n$ reaches only after at least $h - 1 - n$ steps. Hence (2) holds at a point of the layer $n$ with $L^{-3(h-1-n)}B_3\varepsilon_U$ at the index $n$, which is at most $L^3B_3\varepsilon_U$ at the index $h$.

\emph{The consequences.} The hypothesis \cite[(14)]{BalabanVAR} asks that the base belong to the space of regular configurations with the parameter $C_1B_3\varepsilon_1$ and that its averages be within $C_1\varepsilon_1$ of the data, with one constant $C_1$; the first holds with $C_1 = L^3$ by what was shown, and the data of the problem of $\mathbf B''_k$ of \cite[(1.54)]{BalabanLFR2} differ from the averages of $U_{1,2}$ on $Z''_k \cap \Omega''^{\sim2}_{h+1}$ only, by $\tilde Q_\cdot(\eta\mathbf A_0)$, which is bounded by \cite[(1.87), p.~197]{BalabanLFR1}. The conditions \cite[(3.35), (3.36)]{BalabanPBG} follow from the membership in the space of regular configurations by \cite[Sect.~F and Proposition~6, pp.~97--99]{BalabanREG}, which use no criticality.
\end{proof}

The collar of $\mathbf B_1$ is the second obstacle, and we remove it by a truncation. The two runs build the collar from the same set $\tilde\Omega$ by the same rule, and they differ only next to its boundary; a walk which comes near the boundary and returns to a point outside the collar has crossed all the layers of the collar twice, and Ba{\l}aban's bound of a walk makes it small by a factor $L^{-4}$ for each layer.

\begin{lemma}[The Collar of the Set $\mathbf B_1$ in the Two Runs]\label{lem:collar-two-runs}
Let $\tilde\Omega = \tilde\Omega^A_0 \supset \tilde\Omega^A_1 \supset \dots \supset \tilde\Omega^A_h$ and $\tilde\Omega = \tilde\Omega^B_0 \supset \dots \supset \tilde\Omega^B_{h+1}$ be the domains of the collars $\mathbf B_h(\tilde\Omega)$ of run~A and $\mathbf B_{h+1}(\tilde\Omega)$ of run~B \cite[p.~256, ll.~57--62]{BalabanCRE}, and let $L > 2d$. Let $\delta_t > 0$ be a constant of $d$ and $L$ such that the terms of the walks of the decoupled operators of this subsection, at the parameters $\abs{s(\Delta)} \leq e^{\kappa_1}$, and the kernels of these operators decay at the rate $2\delta_t$ at least \cite[(1.7), (1.11), (1.21)]{BalabanRG2}, \cite[(3.42), (3.132), (3.133)]{BalabanPBG}, and let
\begin{equation}\label{eq:truncation-condition}
    \delta_tLM_1 \geq 4\log L\;,
\end{equation}
a lower bound on $M_1$ of the kind of the condition (J-W) of Lemma~\ref{lem:joins-block-size}.
\begin{enumerate}[label=(\alph*)]
\item $\tilde\Omega^A_n = \tilde\Omega^B_{n+1}$ for every $n \geq 2$. The two collars therefore have the same layers at the shared scales $n \geq 2$, and they differ only on the strip $\tilde\Omega \setminus \tilde\Omega^A_2$, which lies within the $\ell^1$-distance $(d + 1)M_1a_2$ of $\partial\tilde\Omega$.
\item For $3 \leq m \leq h$ let $\mathfrak F_m \coloneqq \tilde\Omega \setminus \tilde\Omega^A_m$ be the union of the layers of the collar below $m$, and for one of the operators $T$ of the problem of $\mathbf B_1$ which have walk expansions, at the parameters $\abs{s(\Delta)} \leq e^{\kappa_1}$, let $T^{>m}$ be the sum of the terms of the walks $\omega$ for which no set $\tilde X^5_i$ meets $\mathfrak F_m$. Then $T^{>m}$ satisfies the bounds of $T$; the walks of $T^{>m}$, their cubes, and their weights are the same in the two runs; and the statements of Paper, Theorem~\ref{P-thm:walk-difference}, of Proposition~\ref{prop:two-lattices}(c), (h), and of Lemmata~\ref{lem:scalar-two-lattices}, \ref{lem:vertex-pi}, \ref{lem:vertex-two}, and~\ref{lem:composite-difference} hold for the operators $T^{>m}$ of the two runs, at an aligned pair which satisfies their hypotheses on $\tilde\Omega^A_{m-1}$, with $j_{\min} = m - 1$.
\item For a source or a datum $f$ and a point $x$ of the determining set, in each run and in each of the norms of \cite[(3.42)--(3.47), (3.132), (3.133)]{BalabanPBG},
\[
    \bigl\lvert\bigl((T - T^{>m})f\bigr)(x)\bigr\rvert \leq C_{\mathrm{tr}}\,e^{-\delta_td(x,\mathfrak F_m)}\,[T]_x\sup\abs f\;,
\]
where $[T]_x$ is the bound of $Tf$ at $x$ for a normalised $f$, $d$ is the distance \cite[(2.46)]{BalabanPRT2} of $\mathbf B_1$, and $C_{\mathrm{tr}}$ is a constant of $d$, $L$, $\kappa_1$, and the constants of \cite{BalabanPBG}. If $x$ lies in a layer $n \geq m$ of the collar, or outside the collar with $n \coloneqq h$, then $e^{-\delta_td(x,\mathfrak F_m)} \leq L^{-4(n-m)}$.
\end{enumerate}
\end{lemma}

\begin{proof}
\emph{(a).} In physical units run~A's $\tilde\Omega^A_n$ and run~B's $\tilde\Omega^B_{n+1}$ are unions of the cubes of the side $c_n \coloneqq M_1a_n$ of one partition, selected by the rule that the $\ell^1$-distance to the complement of the preceding domain is at least $c_n$ \cite[p.~256, ll.~57--62]{BalabanCRE}, \cite[p.~223, l.~16]{BalabanPRT2}; run~B has in addition the domain $\tilde\Omega^B_1$, built with the cubes of the side $c_0$. The partitions are nested, so that the distance of a cube of the side $c_n$ from $\tilde\Omega^c$ is a multiple of $c_n$ up to $d$ lattice spacings. The complements of $\tilde\Omega^A_1$ and of $\tilde\Omega^B_2$ both contain $\tilde\Omega^c$ and lie within the $\ell^1$-distance $2dc_1$ of it, each removing at most two layers of cubes of the side $c_1$. Let $Q$ be a cube of the side $c_2$. If $Q$ touches $\tilde\Omega^c$, or if its distance from $\tilde\Omega^c$ is $c_2$, it is selected in neither run, a removed layer of cubes of the side $c_1$ lying between it and $\tilde\Omega^c$ in the second case. If the distance is at least $2c_2$, it is selected in both, because $2c_2 - 2dc_1 \geq c_2$ for $L > 2d$. Hence $\tilde\Omega^A_2 = \tilde\Omega^B_3$, and the rule gives $\tilde\Omega^A_n = \tilde\Omega^B_{n+1}$ for $n \geq 2$ by induction. The strip lies within $c_2 + 2dc_1 \leq (d + 1)c_2$ of $\partial\tilde\Omega$.

\emph{(b).} The bounds of $T$ are sums of the bounds of the terms of its walks \cite[(3.107), (3.108)]{BalabanPBG}, \cite[(1.7), (1.11)]{BalabanRG2}, so that every partial sum satisfies them. A term of a walk is a product of local operators of its cubes, of partition functions, and of averages, and the local operators of a cube $\square$ of a layer $n$ are those of the sequence $\Omega_{n+1}(\square) = \tilde\square^3 \cap B^{n+1}(\Lambda_{n+1})$, $\Omega_n(\square) = \tilde\square^4$, and cubes about the centre of $\square$ below, with $\Omega_0(\square) \subset \tilde\square^5$ \cite[p.~408, l.~29--p.~409, l.~3]{BalabanPBG}; they depend on the determining set only through $\square$ and the trace of the next layer on $\tilde\square^3$, and on the configuration on $\tilde\square^5$. At the block size of Lemma~\ref{lem:joins-block-size} a set $\tilde X^5$ of a cube of a layer $n$ meets at most the layers $n - 1$, $n$, and $n + 1$. A walk no set $\tilde X^5_i$ of which meets $\mathfrak F_m$ has therefore its cubes in the layers at least $m$, their enlargements in the layers at least $m - 1 \geq 2$, and by (a) these cubes, the traces of the layers on them, and the partition functions, as functions on the torus, are the same in the two runs, the local sequences of run~B having one finer shell as on a set at shared scales (Paper, Section~\ref{P-sec:walk-expansions}). The proofs of the results named in (b) compare the two runs walk by walk and factor by factor, with the defect of a local operator at the scale of its cube (Paper, proof of Theorem~\ref{P-thm:walk-difference}, ``The factors''; proofs of Lemmata~\ref{lem:scalar-two-lattices} and~\ref{lem:composite-difference}, ``The decoupled operators''), and they use of the set nothing else; the vertices \eqref{eq:vertex-pi} and \eqref{eq:vertex-two} are compared between operators of the kinds already compared. They apply to the walks of $T^{>m}$ without change.

\emph{(c).} A walk of $T - T^{>m}$ has a set $\tilde X^5_i$ which meets $\mathfrak F_m$, so that the quantity $d(\omega, y, y')$ of Paper, \eqref{P-eq:walk-bound}, for the block $y$ of $x$, is at least $d(x, \mathfrak F_m)$ less the weighted diameter of one enlarged cube, a constant of $d$ and of the block size. At the parameters $\abs{s(\Delta)} \leq e^{\kappa_1}$ the terms of the walks decay at the rate $2\delta_t$ at least; one half of it gives the factor $e^{-\delta_td(x,\mathfrak F_m)}$, and with the other the sum over the walks converges as in \cite[Theorem~3.7]{BalabanPBG}. For the last statement, by \eqref{eq:truncation-condition} it suffices that $d(x, \mathfrak F_m) \geq (n - m)LM_1$: a contour from $\mathfrak F_m$ to $x$ crosses the layers $m, \dots, n - 1$ of the collar, each of which is at least $LM_1$ wide in its own units (Lemma~\ref{lem:layer-widths}), and a contour which crosses such a layer has in it the length at least $LM_1$ in the weights of \cite[(2.46)]{BalabanPRT2}.
\end{proof}

A function of the form \cite[(179), (180)]{BalabanVAR} is determined by its operators $H$, $H_0$, $\tilde G$ and by its data through two fixed-point equations, and the next lemma says that it inherits from them the decay of a change. We state it in the form in which it is used: two functions whose operators differ far from a point differ there by little.

\begin{lemma}[A Function and Its Truncation]\label{lem:function-truncation}
In one run, let $\mathbf H$ and $\mathbf H^>$ be the functions \cite[(1.2)--(1.6)]{BalabanRG2} formed, at one pair and at the same data $B'$, with two families $(H, H_0, \tilde G)$ and $(H^>, H^>_0, \tilde G^>)$ of weighted sums of walks of one determining set, both with the bounds \cite[(1.18)]{BalabanRG2} of the decoupled operators, such that for a function $\varpi$ on the set with $0 < \varpi \leq 1$ and $\varpi(y) \leq e^{\delta_td(x,y)}\varpi(x)$, and a constant $C_{\mathrm{tr}}$,
\begin{equation}\label{eq:operator-truncation}
    \bigl\lvert\bigl((T - T^>)f\bigr)(x)\bigr\rvert \leq C_{\mathrm{tr}}\,\varpi(x)\,[T]_x\sup\abs f
\end{equation}
for each of the three operators and each of their norms. There is a constant $C_\varpi$ of $C_{\mathrm{tr}}$, $d$, $L$, $\kappa_1$, and the constants of \cite{BalabanPBG, BalabanVAR} such that, for $\varrho = 4B_Te^{2\cdot6^d\kappa_1}\abs{B'}$ small, $N(\mathbf H - \mathbf H^>; x) \leq C_\varpi\varpi(x)\varrho$ at every point $x$. The same holds for the functions of the form \cite[(111), (112)]{BalabanVAR}, with the operators $(H, H_1, \mathfrak G)$ and with the current of the background in the data.
\end{lemma}

\begin{proof}
We write $N_\varpi(X) \coloneqq \sup_x\varpi(x)^{-1}N(X; x)$. An operator $K$ with a kernel which decays at a rate of at least $2\delta_t$ maps a field $X$ to a field with $N(KX; x) \leq C\varpi(x)N_\varpi(X)$, because $\varpi(y) \leq e^{\delta_td(x,y)}\varpi(x)$ and the sum of $e^{-\delta_td(x,y)}$ over the blocks converges \cite[(2.58)]{BalabanPRT2}; the operators of the two families have such kernels \cite[(3.42), (3.132), (3.133)]{BalabanPBG}. The equations are $\mathbf H = A' - HD(A')$, $A' = H_0B' + \mathscr A_0$, $\mathscr A_0 = -\tilde G[\Delta^{(2)}H_0B' + (\delta V/\delta A')(A')]$ \cite[(179), (180)]{BalabanVAR}, and the same with the second family. Subtracting, the differences $\delta A' \coloneqq A' - A'^>$ and $\delta\mathscr A_0$ satisfy
\[
    \delta A' = (H_0 - H^>_0)B' + \delta\mathscr A_0\;,
    \qquad
    \delta\mathscr A_0 = -(\tilde G - \tilde G^>)[\,\cdots] - \tilde G^>\Bigl[\Delta^{(2)}(H_0 - H^>_0)B' + \Bigl(\frac{\delta V}{\delta A'}\Bigr)(A') - \Bigl(\frac{\delta V^>}{\delta A'}\Bigr)(A'^>)\Bigr]\;,
\]
where $V^>$ is the functional \cite[(80)]{BalabanVAR} formed with $H^>$. The terms with the differences of the operators are bounded by \eqref{eq:operator-truncation}, applied to arguments of the sizes $\abs{B'}$ and $C_4\varrho^2$, by $C\varpi(x)\varrho$. For the last bracket we use the form of the functional derivative, \cite[(85)--(96)]{BalabanVAR}: each of its terms at a bond is a product of values, at this bond and at the plaquettes and the block which contain it, of the field $A'$, of its first differences, of the datum $D(A')$, and of the field $HD(A')$, with at most two kernels, of $H$ and of the inverse of the block form, between a factor and the current or a datum at another point. A difference of two such terms, at two fields of the size $\varrho$ and with the two operators $H$ and $H^>$, is therefore at most $O(1)C_4\varrho$ times the sum over the blocks $y$, with the decay of these kernels, of $N(\delta A'; y) + \varpi(y)\varrho$, in the norm $\abs\cdot_{(-3)}$; this is the proof of \cite[Proposition~4]{BalabanVAR} and of \cite[(186), (187)]{BalabanVAR}, read term by term. Hence the bracket is at most $O(1)C_4\varrho\,\varpi(x)(N_\varpi(\delta A') + \varrho)$ at $x$, and with the bound $B_T$ of $\tilde G^>$
\[
    N_\varpi(\delta A') \leq C\varrho + O(1)B_TC_4\varrho\bigl(N_\varpi(\delta A') + \varrho\bigr)\;,
\]
so that $N_\varpi(\delta A') \leq 2C\varrho$ once $O(1)B_TC_4\varrho \leq \frac{1}{4}$; the left side is finite because the lattice is finite and $\varpi > 0$. Finally $\mathbf H - \mathbf H^> = \delta A' - (H - H^>)D(A') - H^>(D(A') - D(A'^>))$, and the datum $D$ is a local function of the second order of its field \cite[(44), (54)]{BalabanVAR}. For the form \cite[(111), (112)]{BalabanVAR} the equations are $A_1 = -\mathfrak GJ - \mathfrak G(\delta V/\delta A')(A_1 + H_1B)$ and $\mathbf H = A_1 + H_1B - HD(A_1 + H_1B)$, and the same argument applies with $\mathfrak G$ in the place of $\tilde G$, $H_1$ in that of $H_0$, and the term $(\mathfrak G - \mathfrak G^>)J$ among the differences of the operators.
\end{proof}

We turn to the function of \cite[(1.54)]{BalabanLFR2}. About a base $U_0$ which satisfies \cite[(14)]{BalabanVAR} without being critical, \cite[Sects.~D, E]{BalabanVAR} represent the minimiser as $(\exp i\eta\mathbf H\,U_0)$ up to a gauge transformation, with
\begin{equation}\label{eq:form-111}
    A_1 + \mathfrak GJ + \mathfrak G\Bigl(\frac{\delta V}{\delta A'}\Bigr)(A_1 + H_1B) = 0\;,
    \qquad
    \mathbf H = A_1 + H_1B - HD(A_1 + H_1B)\;,
\end{equation}
where $J$ is the current of $U_0$, $B$ the datum, $V$ the functional \cite[(80)]{BalabanVAR}, and $H_1$, $\mathfrak G = G_1\mathfrak P^*$ the operators \cite[(3.129), (3.148)]{BalabanPBG} \cite[(111), (112), p.~294; Proposition~5]{BalabanVAR}. The operators of this form are built with the series
\begin{equation}\label{eq:series-g1}
    G_1 = \sum_{n \geq 0}G_a\bigl((\Delta'_\pi + \Delta^{(2)}_\pi)G_a\bigr)^n\;,
    \qquad
    \Delta^{(2)}_\pi = P_\pi^*\Delta^{(2)}P_\pi\;,
    \qquad
    P_\pi = 1 - DG'RD^*\;,
\end{equation}
of \cite[(3.135), (3.138)]{BalabanPBG}, in which the vertex $\Delta^{(2)}$ stands between two projections, and $\mathfrak G = G_1 - G_1DRD^*G_1 - G_1Q^*(QG_1Q^*)^{-1}QG_1$ \cite[(3.153), p.~426]{BalabanPBG}. Lemma~\ref{lem:composite-difference} compares the series with the vertex $\Delta'_\pi$ and the series with the vertex $\Delta^{(2)}$; the next lemma adds the projected vertex and the middle term of $\mathfrak G$.

\begin{lemma}[The Operators of the Form (111) on Two Lattices]\label{lem:polarised-operators}
Under the hypotheses of Lemma~\ref{lem:composite-difference}, and for every $\beta_c \in (0, \frac{1}{60}]$, each of the differences $G^B_1 - \mathcal LG^A_1\mathcal L^*$, $H^B_1 - \mathcal LH^A_1$, $\mathfrak G^B - \mathcal L\mathfrak G^A\mathcal L^*$, and of the averaged differences $Q_1H^B_1 - H^A_1$, $Q_1\mathfrak G^B - \mathfrak G^A\mathcal L^*$, is at most $C'_{\mathfrak d}\mathfrak d_H$ in the supremum and $C'_{\mathfrak d}L^{\beta_cj}\mathfrak d_H$ in the first differences, at the points of a layer $j$ and relative to the bounds of the operators, with a constant $C'_{\mathfrak d}$ which depends on $\beta_c$ and on the constants of that lemma. The same holds walk by walk for the decoupled operators.
\end{lemma}

\begin{proof}
\emph{The projected vertex as a source.} For a field $A$ of either run let $\omega_A \coloneqq G'RD^*A$, so that $P_\pi A = A - D\omega_A$, and let $g_A \coloneqq \Delta^{(2)}P_\pi A$ and $\varphi'_A \coloneqq RG'D^*g_A$. Then $\Delta^{(2)}_\pi A = P_\pi^*g_A = g_A - D\varphi'_A$. With $N \coloneqq \sup(a\abs A + a^2\abs{\nabla A})$ over the blocks, weighted by the decay of the kernels, $\abs{D\omega_A} \leq c_\omega a^{-1}N$ \cite[(3.42), (3.49)]{BalabanPBG}, hence $\abs{g_A} \leq C\alpha_Ja^{-3}N$ \cite[(3.137)]{BalabanPBG}, and $\varphi'_A$ is H\"older continuous of the exponent $\beta^+$ with the constant $Ca^{1-\beta^+}\sup\abs{g_A}$, by the second entry of \cite[(3.43)]{BalabanPBG}. The source $(\Delta'_\pi + \Delta^{(2)}_\pi)A$ is therefore, as the source $\Delta'_\pi A$ in the proof of Lemma~\ref{lem:composite-difference}(b), a bounded source plus the gradient of a function which is H\"older continuous of the exponent $\beta^+$.

\emph{The two runs at the vertex.} Let $A$ be a field of run~A with $N_\beta(A; \cdot)$ finite, and form run~B's objects at $\mathcal LA$ and at $J^B = Q^{\mathrm{pl}*}_1J^A$. For a field $Y$ and a field $X$ of run~B the composition of the averages gives $C^{(2)B}(Y, X) = C^{(2)A}(Q_1Y, Q_1X) + \hat Q^Ac^{(2)}_1(Y, X)$ for the bilinear parts, and with $\langle Y, \Delta^{(2)}X\rangle = 2\langle HC^{(2)}(Y, X), J\rangle$ and $Q^{\mathrm{pl}}_1 = Q_1 - \bar D\Lambda_1$,
\[
    \tfrac{1}{2}\langle Y, \Delta^{(2)B}X\rangle_B - \tfrac{1}{2}\langle Q_1Y, \Delta^{(2)A}Q_1X\rangle_A = \bigl\langle(Q_1H^B - H^A)C^{(2)A}(Q_1Y, Q_1X) + Q_1H^B\hat Q^Ac^{(2)}_1(Y, X) - \bar D\Lambda_1H^BC^{(2)B}(Y, X),\, J^A\bigr\rangle_A\;,
\]
which is the identity before Lemma~\ref{lem:vertex-two} with a general first argument. Its three terms are bounded as in the proof of that lemma, by $C\alpha_J\mathfrak d_H\,a^{-3}$ times the norm of $Y$ in $L^1$ of a block and the supremum of $a\abs X$, with the decay of the kernel of $H$. We take $X = P^B_\pi\mathcal LA$. By the identity $Q_1D_B = D_AQ'_1$ of \eqref{eq:plain-dressed} and $Q_1\mathcal L = 1$, $Q_1P^B_\pi\mathcal LA = A - D_AQ'_1\omega^B = P^A_\pi A - D_Ae$ with $e \coloneqq Q'_1\omega^B - \omega^A$, and $\abs{D_Ae} \leq C\mathfrak s_*^{\beta/(1+\beta)}a^{-1}N_\beta$ (proof of Lemma~\ref{lem:vertex-pi}, ``The fields of the second argument''). Since the form $\Delta^{(2)A}$ is bounded on the pairs of an integrable and a bounded field, we obtain, in the supremum,
\begin{equation}\label{eq:projected-source}
    \bigl\lvert g^B - Q_1^*g^A\bigr\rvert \leq C\alpha_J\mathfrak d_H\,a^{-3}N_\beta\;.
\end{equation}
For the functions $\varphi'$ we write $\varphi'^B - \mathcal L'\varphi'^A = R^BG'^BD_B^*(g^B - Q_1^*g^A) + [R^BG'^BD_B^*Q_1^* - \mathcal L'R^AG'^AD_A^*]g^A$. The first term is bounded by \eqref{eq:projected-source} and the entry $G'\nabla^*$ of \cite[(3.42)]{BalabanPBG}. The second is the left side of \eqref{eq:rough-divergence} with $Q_1^*$ in the place of $Q^{\mathrm{pl}*}_1$, and the proof of Lemma~\ref{lem:rough-divergence} applies to it: the field $Q_1^*\mathfrak m$ is bounded by Lemma~\ref{lem:local-relations}(b), and in the computation of the local means the identity $Q_1D_B\Psi^B = D_AQ'_1\Psi^B$ takes the place of Lemma~\ref{lem:local-relations}(a), so that the term with $\bar D\Lambda_1$ does not occur. Hence $\abs{\varphi'^B - \mathcal L'\varphi'^A} \leq C\alpha_J\mathfrak d_H\,a^{-2}N_\beta$.

\emph{The form.} For fields $A_1$, $A_2$ of run~A, moving the projection of the first argument to the source,
\[
    \langle\mathcal LA_1, \Delta^{(2)B}_\pi\mathcal LA_2\rangle_B - \langle A_1, \Delta^{(2)A}_\pi A_2\rangle_A = \bigl\langle\mathcal LA_1, g^B - Q_1^*g^A\bigr\rangle_B - \bigl[\langle D_B^*\mathcal LA_1, \varphi'^B\rangle_B - \langle D_A^*A_1, \varphi'^A\rangle_A\bigr]\;,
\]
because $\langle\mathcal LA_1, Q_1^*g^A\rangle_B = \langle A_1, g^A\rangle_A$. The bracket is treated as the second, third, and fourth terms of \eqref{eq:vertex-difference}, with $\varphi'$ in the place of $\varphi$: by Lemma~\ref{lem:local-relations}(a), (b) it is $\langle D_B^*\mathcal LA_1, \varphi'^B - \mathcal L'\varphi'^A\rangle_B - \langle\mathfrak e(\mathcal LA_1), D_A\varphi'^A\rangle_A - \langle\mathfrak e_0(\mathcal LA_1), \varphi'^A\rangle_A$, and the last two terms are bounded as there, the first differences of $\varphi'^A$ by \cite[(3.44)]{BalabanPBG} and \eqref{eq:crude-holder}. The projected vertex therefore satisfies \eqref{eq:vertex-pi}, with a larger constant.

\emph{The series $G_1$ and the operator $H_1$.} The telescoping \eqref{eq:series-telescoping} with the vertex $\Delta'_\pi + \Delta^{(2)}_\pi$ in the place of $\Delta'_\pi$ has the three kinds of terms of the proof of Lemma~\ref{lem:composite-difference}(b), and they are bounded as there: the difference of the propagators on a bounded source by Proposition~\ref{prop:two-lattices}(h), on the source of the vertex by Lemma~\ref{lem:composite-difference}(a) and the first step, and the difference of the vertices by Lemma~\ref{lem:vertex-pi} and the third step; the factors on the left preserve the smallness with the factor $C\alpha_J$ each. For $H_1 = G_1Q^*(QG_1Q^*)^{-1}$ the resolvent identity of that proof holds with $G_1$ in the place of $G_\pi$.

\emph{The operator $\mathfrak G$.} With $F(X) \coloneqq X - XQ^*(QXQ^*)^{-1}QX$ we have $\mathfrak G = F(G_1) - G_1DRD^*G_1$, and $F(G^B_1) - F(\mathcal LG^A_1\mathcal L^*)$ is bounded as for $\tilde G$ in Lemma~\ref{lem:composite-difference}(b). For the middle term let $f$ be a bounded source of run~B, $u^B \coloneqq G^B_1f$, $u^A \coloneqq G^A_1\mathcal L^*f$, and $\psi^X \coloneqq R^XD_X^*u^X$; the term is $G_1D\psi$ in either run. The first differences of $u^X$ are H\"older continuous \cite[(3.43)]{BalabanPBG}, so that $\psi^X$ is H\"older continuous of the exponent $\beta^+$ and $G_1D\psi$ is a propagator on a gradient source, to which Lemma~\ref{lem:composite-difference}(a) and the series apply. By Lemma~\ref{lem:local-relations}(b),
\[
    Q'_1D_B^*u^B - D_A^*u^A = D_A^*\bigl(Q_1u^B - u^A\bigr) - D_A^*\mathfrak e(u^B) - \mathfrak e_0(u^B)\;,
\]
in which $Q_1u^B - u^A = (Q_1G^B_1 - G^A_1\mathcal L^*)f$ is small with its first differences by what has been proved, and the other two terms carry the factor $(\varepsilon/a)^\beta$ by the H\"older continuity of $\nabla u^B$; and $Q'_1R^B - R^AQ'_1 = -(Q'_1P^B - P^AQ'_1)$ is given by the identity for the projection in the proof of Lemma~\ref{lem:scalar-two-lattices}, each term of which contains the defect of $G'$ once. Hence $\abs{Q'_1\psi^B - \psi^A}$ is at most $CL^{\beta_cj}\mathfrak d_H$ times the size of $\psi$, and $\psi^B - \mathcal L'\psi^A = (1 - \mathcal L'Q'_1)\psi^B + \mathcal L'(Q'_1\psi^B - \psi^A)$ is small in the supremum, the first term by the H\"older continuity of $\psi^B$. Now
\[
    G^B_1D_B\psi^B - \mathcal LG^A_1D_A\psi^A = \bigl(G^B_1 - \mathcal LG^A_1\mathcal L^*\bigr)D_B\psi^B + \mathcal LG^A_1\bigl(\mathcal L^*D_B\psi^B - D_A\psi^A\bigr)\;,
\]
and at a bond $x$ of run~A the second term is $\langle D_B^*\mathcal LA_1, \psi^B\rangle_B - \langle D_A^*A_1, \psi^A\rangle_A$ with $A_1 = G^A_1(\cdot, x)$, which is bounded as the bracket of the third step. For the first differences in $x$ the kernel $\nabla G_1\nabla^*$ costs the factor $L^{\beta_cj}$ as in Lemma~\ref{lem:composite-difference}, and the factor of $Q'_1\psi^B - \psi^A$ makes it $L^{2\beta_cj}$; applying the argument with $\frac{1}{2}\beta_c$ gives the statement. The averaged differences follow with $Q_1\mathcal L = 1$. For the decoupled operators every step bounds the terms of one composite walk, as in the proofs of Lemmata~\ref{lem:scalar-two-lattices}, \ref{lem:vertex-pi}, and~\ref{lem:composite-difference}, the projection $R$ being expanded through \cite[(3.90), (3.98)]{BalabanPBG}.
\end{proof}

\begin{lemma}[The Function of the Form (111) on Two Lattices]\label{lem:field-difference-111}
Let the hypotheses of Lemma~\ref{lem:field-difference} hold on a set $Y$ whose layers lie at the shared scales between $j_{\min} \geq m_{\mathrm{rc}}$ and $j_{\max}$, at a pair $(U^A, J^A)$ and its lift. Let $\mathbf H^A$ and $\mathbf H^B$ be the functions of the two runs defined by
\[
    A_1 + \mathfrak s + \mathfrak G\Bigl(\frac{\delta V}{\delta A'}\Bigr)(A_1 + H_1B) = 0\;,
    \qquad
    \mathbf H = A_1 + H_1B - HD(A_1 + H_1B)\;,
\]
with the decoupled operators at the same parameters, with data $B^X$ and fields $\mathfrak s^X$ such that the fields $\mathfrak s^X$, $H^X_1B^X$, $A^X_1$, and $\mathbf H^X$ are at most $\varrho$ in the norm $N$, $\varrho$ being as small as the contraction \cite[(116)--(121)]{BalabanVAR} requires. Then, at the points whose complex coordinates lie within a fifth of the radii,
\begin{equation}\label{eq:field-difference-111}
    N(\mathbf H^B - \mathcal L_\iota\mathbf H^A) \leq C_{111}\Bigl[L^{\beta_cj_{\max}}\bigl(C_{\mathfrak d}L^{-\gamma_Hj_{\min}} + L^{-j_{\min}}\bigr)\varrho + N(\mathfrak s^B - \mathcal L_\iota\mathfrak s^A) + \sup\abs{B^B - B^A}\Bigr]\;,
\end{equation}
with a constant $C_{111}$ independent of $k$, $K$, $\varepsilon$, the coupling, and the torus. The bound \eqref{eq:field-difference} of Lemma~\ref{lem:field-difference} holds in the same form, with the first term of the bracket and the last, for the functions \cite[(179), (180)]{BalabanVAR} at data $B'^X$ on such a set.
\end{lemma}

\begin{proof}
The proof is that of Lemma~\ref{lem:field-difference}, and we name the changes. That proof uses the hypothesis on the layers twice, through $\varepsilon/a_j \leq L^{-j_{\min}}$ and through the number $\ell_c = \sup_jL^{\beta_cj/2}$, which is at most $L^{\beta_cj_{\max}/2}$ here; its result is $N(\Delta_{\mathbf H}) \leq C\ell_c^2\mathfrak d_0\varrho$, which is the first term of the bracket. The datum enters it through the field $H_0B'$ alone, linearly, and two data give the further term $H^B_0(B'^B - B'^A)$, bounded by the last term of the bracket; this proves the last statement.

For the form \eqref{eq:form-111} the triple $(A_1, \mathfrak G, H_1)$ takes the place of $(\mathscr A_0, \tilde G, H_0)$. The identity \eqref{eq:field-difference-identity} holds with $H_1B$ in the place of $H_0B'$ and $\Delta_1 \coloneqq A^B_1 - \mathcal L_\iota A^A_1$ in that of $\Delta_0$, and subtracting the two fixed-point equations,
\[
    \Delta_1 = -(\mathfrak s^B - \mathcal L_\iota\mathfrak s^A) - \mathfrak G^B\bigl[V'^B(A'^B) - V'^B(\mathcal L_\iota A'^A)\bigr] - (\mathfrak G^B - \mathcal L_\iota\mathfrak G^A\mathcal L_\iota^*)V'^B(\mathcal L_\iota A'^A) - \mathcal L_\iota\mathfrak G^A\sigma\;,
\]
with $A' = A_1 + H_1B$ and $\sigma = \mathcal L_\iota^*V'^B(\mathcal L_\iota A'^A) - V'^A(A'^A)$; there is no term $\mathsf T$, the form $\Delta^{(2)}$ being contained in the operators. The second term is the contraction \cite[(119)--(121)]{BalabanVAR}, the third is Lemma~\ref{lem:polarised-operators} on a source of the size $C_4\varrho^2$, and the operators $H_1$ and $H$ in the other terms are compared by the same lemma and by Lemma~\ref{lem:composite-difference}. One point is new. The functional of \eqref{eq:form-111} is that of \cite[(80)]{BalabanVAR} with the form $\Delta_\pi$, whereas the identity \eqref{eq:functional-difference} is derived for the functional with the form $\Delta$; their difference is $\langle A', \Delta'_\pi HD(A')\rangle - \frac{1}{2}\langle HD(A'), \Delta'_\pi HD(A')\rangle$. Accordingly $\sigma$ is the source of the proof of Lemma~\ref{lem:field-difference} plus the two-run difference of the gradient of this expression, which is the difference of the forms $\Delta'_\pi$ of the two runs at the arguments $A'$ and $HD(A')$ and their derivatives; by Lemma~\ref{lem:vertex-pi}, with the comparison of $H$ and of the data $D$ in the first two paragraphs of that proof, it is at most $C\alpha_J\mathfrak d_H\varrho^2$ in the norm in which $\mathfrak G^A$ is applied to it. All the other terms of $\sigma$ are bounded as there, the artefact included, with the kernels of $\mathfrak G^A$ and of $\nabla\mathfrak G^A$ in the places of those of $\tilde G^A$; their norms in $L^1$ are bounded by \cite[Theorem~3.13]{BalabanPBG}. The two contractions close as there, and $\mathbf H^B - \mathcal L_\iota\mathbf H^A$ is obtained from $\Delta_1$ by \eqref{eq:first-contraction}.
\end{proof}

In the localisation of \cite[(1.61)]{BalabanLFR2} the term $\mathfrak GJ$ of \eqref{eq:form-111} does not stand as it is. Ba{\l}aban writes it, by the criticality of the current $J_1$ of $U_1$ for the set $\mathbf B_1$ and by $\mathbf G''_kQ^*_{\mathbf B''_k} = 0$, as
\begin{equation}\label{eq:source-155}
    \mathfrak s \coloneqq \mathbf G''_k(1 - \theta)Q^*_{\mathbf B_1}\theta_1H^*_{\mathbf B_1}J_1 + \mathbf G''_k\theta J_{1,2}\;,
\end{equation}
with a function $\theta \in C^\infty_0((\Omega''_{h+1})^c)$ which is one on $\tilde\Omega^c$ and on the boundary layer of the width $2M_1$ at the scale $h$ and changes to zero near the inner boundary of this layer, and a function $\theta_1$ which is one where $\theta \neq 0, 1$ and changes to zero on a neighbourhood of that set; ``we use this fact to localize the considered expression'' \cite[(1.55), p.~371]{BalabanLFR2}. The function of the construction is therefore the function of Lemma~\ref{lem:field-difference-111} with the field \eqref{eq:source-155}, in each run, and at a real configuration of the run it is the function of \eqref{eq:form-111}. Ba{\l}aban fixes $\theta$ and $\theta_1$ by these properties only. The sets in them are sets of the scale $h$, the same in the two runs, and we take in both runs the same two functions of the points of the torus, with $a_h\abs{\partial\theta} + a_h^2\abs{\partial\partial\theta} \leq c_\theta$ and the same for $\theta_1$, and with the support of $\theta_1$ in the layers of the collar of the indices at least $h - 1$ and in $\tilde\Omega^c$.

\begin{lemma}[The Source of (1.55) in the Two Runs]\label{lem:source-155}
Let \eqref{eq:truncation-condition} hold, let $3 \leq m \leq h - 1$, and let the fields \eqref{eq:source-155} of the two runs be formed at a real point, run~A's at the pair $(U^A_{1,2}, J^A_{1,2})$ with $J^A_{1,2}$ the current of $U^A_{1,2}$ and run~B's at its lift, with the decoupled operators at the same parameters. Let $\varrho_{\mathfrak s}$ be the supremum of $a^3\abs{J_{1,2}}$ over the layers. Then, with a constant $C_{\mathfrak s}$ independent of $k$, $K$, $\varepsilon$, the coupling, and the torus,
\begin{equation}\label{eq:source-difference}
    N(\mathfrak s^B - \mathcal L_\iota\mathfrak s^A) \leq C_{\mathfrak s}\Bigl[L^{\beta_ck}\bigl(C_{\mathfrak d}L^{-\gamma_H(m-1)} + L^{-(m-1)}\bigr) + L^{-4(h-1-m)}\Bigr]\varrho_{\mathfrak s}\;.
\end{equation}
\end{lemma}

\begin{proof}
Let $\lambda^X \coloneqq H^{X*}_{\mathbf B_1}J^X_1$, a datum on the blocks of $\mathbf B_1$ in the support of $\theta_1$, and $\sigma^X \coloneqq (1 - \theta)Q^{X*}_{\mathbf B_1}\theta_1\lambda^X + \theta J^X_{1,2}$, so that $\mathfrak s^X = \mathbf G''^X_k\sigma^X$ and
\[
    \mathfrak s^B - \mathcal L_\iota\mathfrak s^A = \bigl(\mathbf G''^B_k - \mathcal L_\iota\mathbf G''^A_k\mathcal L_\iota^*\bigr)\sigma^B + \mathcal L_\iota\mathbf G''^A_k\bigl(\mathcal L_\iota^*\sigma^B - \sigma^A\bigr)\;.
\]
The operator $\mathbf G''_k$ is the operator $\mathfrak G$ of the set $\mathbf B''_k$ \cite[p.~370, l.~49]{BalabanLFR2}, whose layers on $Z$ lie at the shared scales between $h$ and $k$, and the source $\sigma^B$ is bounded by a constant times $\varrho_{\mathfrak s}a^{-3}$, the datum $\lambda$ being bounded by \cite[(3.133)]{BalabanPBG} transposed; Lemma~\ref{lem:polarised-operators} bounds the first term by the first term of \eqref{eq:source-difference} at $m - 1 \leq h$.

For the second term we use three facts. First, at the lift the current of run~B is $J^B = Q^{\mathrm{pl}*}_1J^A$, and $Q^{\mathrm{pl}}_1\mathcal L_\iota = 1 - \bar D\Lambda_1\mathcal L_\iota$ by \eqref{eq:plain-dressed} and $Q_1\mathcal L_\iota = 1$; since the covariant divergence of the current of a configuration vanishes, $\langle\bar D\mu, J^A_{1,2}\rangle_A = 0$ for every function $\mu$ \cite[(3.117)]{BalabanPBG}, and therefore $\langle\mathcal L_\iota v, J^B\rangle_B = \langle v, J^A\rangle_A$ for every field $v$ of run~A. The same holds for $J_1$ on $\tilde\Omega$, because $J_{1,2}$ is $J_1$ there and vanishes at the cut (Lemma~\ref{lem:glued-configuration}). Second, the kernel of $\mathcal L_\iota$ decays at the scale $\varepsilon$ \cite[(3.133)]{BalabanPBG}, and the functions $\theta$, $\theta_1$ vary at the scale $a_h$, so that the commutators of $\mathcal L_\iota^*$ with the multiplications by them are at most a constant times $c_\theta\varepsilon/a_h$ times the bound of $\mathcal L_\iota^*$. Third, $\mathcal L_\iota^*Q^{B*} = Q^{A*}$ on the data of the blocks at the shared scales (proof of Lemma~\ref{lem:composite-difference}(b)). With these,
\[
    \mathcal L_\iota^*\sigma^B - \sigma^A = (1 - \theta)Q^{A*}_{\mathbf B_1}\theta_1\bigl(\lambda^B - \lambda^A\bigr) + \mathsf c\;,
\]
where $\mathsf c$ consists of the two commutators and, in the pairing with a field of run~A, is at most a constant times $(\varepsilon/a_h)\varrho_{\mathfrak s}$ times its norm in $L^1$; the term $\theta(\mathcal L_\iota^*J^B - J^A)$ vanishes in these pairings by the first fact. At a block bond $c$ in the support of $\theta_1$, with the unit datum $e_c$, the first fact gives
\[
    \lambda^B(c) - \lambda^A(c) = \bigl\langle H^B_{\mathbf B_1}e_c, J^B_1\bigr\rangle_B - \bigl\langle H^A_{\mathbf B_1}e_c, J^A_1\bigr\rangle_A = \bigl\langle\bigl(H^B_{\mathbf B_1} - \mathcal L_\iota H^A_{\mathbf B_1}\bigr)e_c, J^B_1\bigr\rangle_B\;.
\]
We split $H_{\mathbf B_1} = H^{>m}_{\mathbf B_1} + (H_{\mathbf B_1} - H^{>m}_{\mathbf B_1})$ in both runs. The difference of the truncated operators is bounded by Lemma~\ref{lem:collar-two-runs}(b) and Lemma~\ref{lem:composite-difference}, with the first term of \eqref{eq:source-difference}, and its pairing with the current is bounded as the datum $\lambda$ itself. A walk of $H_{\mathbf B_1} - H^{>m}_{\mathbf B_1}$ from the block of $c$ passes through $\mathfrak F_m$, and $c$ lies in a layer of the index at least $h - 1$; by the proof of Lemma~\ref{lem:collar-two-runs}(c) its term carries, at every point, the factor $e^{-\delta_td(c,\mathfrak F_m)} \leq L^{-4(h-1-m)}$, in each run. Hence $\abs{\lambda^B - \lambda^A}$ is at most a constant times the bracket of \eqref{eq:source-difference} times $\varrho_{\mathfrak s}$, in the units of the data. Finally $\mathcal L_\iota\mathbf G''^A_k$ applied to these sources is bounded in the norm $N$ by \cite[(3.42)]{BalabanPBG} and \eqref{eq:lift-regularity}, and $\varepsilon/a_h \leq L^{-(m-1)}$.
\end{proof}

We can now state the result. On a young matched pair with the finest layer $\nu$, every layer of the sets $\mathbf B_k$ and $\mathbf B''_k$ of an $\mathbf R$-operation of a step $k$ lies at a shared scale between $\nu$ and $k \leq 2\nu$ (Section~\ref{sec:chain-members}), and $h \geq \nu$.

\begin{proposition}[The Localisation Functions of an $\mathbf R$-Operation on Two Fine Lattices]\label{prop:localisation-functions}
Let the hypotheses of Lemmata~\ref{lem:field-difference} and~\ref{lem:composite-difference} hold, with $\gamma_H \geq \frac{1}{12}$ and $\beta_c \leq \frac{1}{60}$, let \eqref{eq:band-condition} and \eqref{eq:truncation-condition} hold with $\delta_t \leq \frac{1}{2}\delta'_0$, let $L > 2d$, and let $\gamma$ be small. There are constants $C_{\mathrm{loc}}$ and $\nu_{\mathrm{loc}}$, independent of $k$, $K$, $\varepsilon$, the coupling, and the torus, with the following property. On a young matched pair with the finest layer $\nu \geq \nu_{\mathrm{loc}}$, at a real point of the region of integration of an $\mathbf R$-operation of a step $k$ and at the points whose complex coordinates lie within a fifth of their radii, let $\mathbf H^X$ be one of the functions $\mathbf H_k$ of \cite[(1.63), (1.64)]{BalabanLFR2}, $\mathbf H_{\mathbf B_1}$ of \cite[(1.62)]{BalabanLFR2}, and $\mathbf H''_k$ of \cite[(1.61)]{BalabanLFR2} of run $X$, at parameters with $\abs{s(\Delta)} \leq \frac{1}{5}e^{\kappa_1}$, the same in the two runs: run~A's at the pair of its base, which is $U_k$, $U_{k,\Lambda}$, or $U_{1,2}$ with its current, and run~B's at the lift of that pair, each with the data of its run. Let $\varrho$ bound the function in the norm $N$ in both runs, and let $d_\Lambda$ be the supremum, over the block bonds of $\Lambda$ at the shared scales, of the logarithms of the ratios of the data $V^{(A)}_\Lambda$ and $M^\cdot(U^{(AL)}_0)$ of the two runs, which the proof of Proposition~\ref{prop:minimisers}(d) bounds. Then, at every point $x$ of the set of the function outside $\Lambda$,
\begin{equation}\label{eq:localisation-functions}
    N(\mathbf H^B - \mathcal L_\iota\mathbf H^A; x) \leq C_{\mathrm{loc}}\bigl(L^{-\nu/26}\varrho + d_\Lambda\bigr)\;,
\end{equation}
and at every site $x$ of run~A outside $\Lambda$ the gauge transformations $u_{k,\Lambda}$, $u_1$, and $u^0_k$ of the three functions satisfy
\begin{equation}\label{eq:localisation-gauge}
    \abs{u^B(x) - u^A(x)} \leq C_{\mathrm{loc}}\bigl(L^{-\nu/26}\varrho + d_\Lambda\bigr)\;.
\end{equation}
\end{proposition}

\begin{proof}
\emph{The function $\mathbf H_k$.} It is the function \cite[(179), (180)]{BalabanVAR} of the set $\mathbf B_k$ about its minimiser $U_k$ at real data, with the datum $-i\log V'_k$ on $\Lambda$, $V'_k = V_k(V^{(A)}_\Lambda)^{-1}$ \cite[(1.58), (1.64)]{BalabanLFR2}. The variables $V_k$ are those of the two runs at equal values, so that the data differ by at most $d_\Lambda$, and Lemma~\ref{lem:field-difference-111}, in its last statement, gives \eqref{eq:localisation-functions} with $j_{\min} = \nu$ and $j_{\max} = k \leq 2\nu$: the first term of the bracket of \eqref{eq:field-difference-111} is at most a constant times $L^{(2\beta_c - \gamma_H)\nu}\varrho \leq L^{-\nu/20}\varrho$.

\emph{The function $\mathbf H_{\mathbf B_1}$.} The configuration $U_{k,\Lambda}$ is a minimiser of $\mathbf B_k$, every block of $\mathbf B_1$ lies in a block of $\mathbf B_k$ of an index at least its own, and by Lemma~\ref{lem:collar-critical} its restriction is the minimiser of $\mathbf B_1$ at its own averages. The representation \cite[(1.57)]{BalabanLFR2} is therefore that of a minimiser of $\mathbf B_1$ about a minimiser of $\mathbf B_1$, and $\mathbf H_{\mathbf B_1}$ is the function \cite[(179), (180)]{BalabanVAR} of this set with the datum $-i\log M^\cdot(U_{k,\Lambda})$ on the blocks of $Z''_k \cap \tilde\Omega$. At the lift this datum is the same in the two runs on the blocks at the shared scales, the average of the lift being the configuration. Let $m \coloneqq h - \lfloor h/8\rfloor$, and let $\mathbf H^{X,>m}$ be the functions formed with the truncated operators of Lemma~\ref{lem:collar-two-runs}(b). Only the data of the blocks of the layers at least $m$ enter them, and by that lemma the proof of Lemma~\ref{lem:field-difference} applies to them with $j_{\min} = m - 1$ and $j_{\max} = k$, which gives, as in the last statement of Lemma~\ref{lem:field-difference-111},
\[
    N(\mathbf H^{B,>m} - \mathcal L_\iota\mathbf H^{A,>m}) \leq C\,L^{2\beta_c\nu}\bigl(C_{\mathfrak d}L^{-\gamma_H(m-1)} + L^{-(m-1)}\bigr)\varrho \leq C'L^{-\nu/26}\varrho\;,
\]
because $m - 1 \geq \frac{7}{8}\nu - 1$, $\gamma_H \geq \frac{1}{12}$, and $2\beta_c \leq \frac{1}{30}$, so that the exponent is at most $-(\frac{7}{96} - \frac{1}{30})\nu + 1$. In each run, by Lemma~\ref{lem:collar-two-runs}(c), the operators and their truncations satisfy \eqref{eq:operator-truncation} with $\varpi(x) \coloneqq e^{-\delta_td(x,\mathfrak F_m)}$, and Lemma~\ref{lem:function-truncation} gives $N(\mathbf H^X - \mathbf H^{X,>m}; x) \leq C_\varpi\varpi(x)\varrho$. A point $x$ outside $\Lambda$ lies outside the collar, which is contained in the band and hence in $\Lambda$ (Lemma~\ref{lem:glued-configuration}), so that $\varpi(x) \leq L^{-4(h-m)} \leq L^4L^{-\nu/2}$; the lift preserves the bound, its kernel decaying at the rate $\delta'_0 \geq 2\delta_t$ in the distance of the fine lattice, which is at least $d$. The three bounds together are \eqref{eq:localisation-functions}.

\emph{The function $\mathbf H''_k$.} Its base $U_{1,2}$ satisfies \cite[(14)]{BalabanVAR} for the set $\mathbf B''_k$ (Lemma~\ref{lem:glued-configuration}), so that \cite[Sects.~D, E]{BalabanVAR} apply and the function is that of \eqref{eq:form-111} with the field \eqref{eq:source-155} and with the datum of \cite[(1.54)]{BalabanLFR2}, the logarithm of the ratio of the averages $M^\cdot(U^{(AL)}_0)$ and of those of $U_{1,2}$ on $Z''_k \cap \Omega''^{\sim2}_{h+1}$, which differs between the two runs by at most $d_\Lambda$ at the lift. The pair of $U_{1,2}$ and its current is regular: by \cite[(1.54)]{BalabanLFR2}, $U_{1,2} = e^{-i\eta\mathbf H''_k}V_0$ with $V_0$ a gauge transform of the real minimiser $U^0_k$ of $\mathbf B''_k$ and with the real field $\mathbf H''_k$, which has the bounds \cite[(19), p.~281]{BalabanVAR} of its values, its first differences, and its Laplacian at the units of the layers, hence H\"older continuous first differences by Lemma~\ref{lem:four-entries}; and its current has a vanishing covariant divergence and is at most $L^3B_3\varepsilon_U$ in the units of the layers of $\mathbf B''_k$, so that it is admissible. The bound \eqref{eq:lift-oscillation} holds at $U_{1,2}$ by the argument for the factors in the proof of Lemma~\ref{lem:field-difference}, with the one factor $e^{-i\eta\mathbf H''_k}$ over the minimiser $V_0$: at a bond of a layer $j$ the oscillation of the first differences of the lifted field is at most $c_e(\beta)(\varepsilon/a_j)^\beta a_j^{-2}$ times the supremum, weighted with the decay at the scale $\varepsilon$, of the H\"older norms of the field in the unit $a_j$ (Lemma~\ref{lem:local-relations}(e)); at the points of the layer $j$ and of the coarser layers these are bounded in their own units, and the points of a finer layer $j' < j$, where they are larger by $L^{(2+\beta)(j-j')}$ at most, lie at a distance of at least $LM_1a_{j-1}$, which the decay at the scale $\varepsilon$ pays as in the case of a minimiser there. Lemma~\ref{lem:field-difference-111} therefore applies on the set $\mathbf B''_k$, with $j_{\min} = \nu$ on $Z^c$ and $j_{\min} = h \geq \nu$ on $Z$, and Lemma~\ref{lem:source-155} with $m = h - \lfloor h/8\rfloor$ bounds $N(\mathfrak s^B - \mathcal L_\iota\mathfrak s^A)$ by a constant times $L^{-\nu/26}\varrho_{\mathfrak s}$, where $\varrho_{\mathfrak s} \leq L^3B_3\varepsilon_U$ is at most a constant times the bound of the function.

\emph{The gauge transformations.} They are the transformations \cite[(106)]{BalabanAVG} of the three functions on the blocks of their sets \cite[p.~373, ll.~29--30, 44; p.~374, ll.~38--39]{BalabanLFR2}, and a block of a point outside $\Lambda$ lies at the shared scales. Lemma~\ref{lem:gauge-transformations-runs} with $d_{\mathbf H}\varrho$ the right side of \eqref{eq:localisation-functions} gives \eqref{eq:localisation-gauge}, since $\varepsilon/a_j \leq L^{-\nu}$.

\emph{The complex points.} Lemmata~\ref{lem:field-difference} and~\ref{lem:field-difference-111} hold at the points whose complex coordinates lie within a fifth of the radii; Lemmata~\ref{lem:collar-two-runs} and~\ref{lem:function-truncation} concern one run and hold on the polydiscs of the parameters and at the complex configurations of the class, at which the operators have their bounds (Lemma~\ref{lem:pbg-complex}); the identities of the proof of Lemma~\ref{lem:source-155} are algebraic, and the current of the lift is defined as $Q^{\mathrm{pl}*}_1J^A$ there; and the gauge transformations are analytic functions of the fields with the Lipschitz constants of the proof of Lemma~\ref{lem:gauge-transformations-runs} on the complex fields of the size $\varrho$.
\end{proof}

% ---------------------------------------------------------------------------------------------------------------
% PART 5 (to follow Parts 2--4): the gauge transformations. Hypothesis on the fields is supplied by Parts 2--4.

\begin{lemma}[The Gauge Transformations of the Two Runs]\label{lem:gauge-transformations-runs}
Let $\mathbf H^A$ and $\mathbf H^B$ be fields of the two runs on a set $Y$ all of whose layers lie at shared scales, over a real pair $(V_0, U^B)$ with $MU^B = V_0$ which satisfies \eqref{eq:minimiser-regularity} with $\varepsilon_1$ at the units of the layers, with $N(\mathbf H^X) \leq \varrho$ and $N(\mathbf H^B - \mathcal L_\iota\mathbf H^A) \leq d_{\mathbf H}\varrho$, and let $u^X$ be the gauge transformation \cite[(106), p.~33]{BalabanAVG} of the configuration $e^{i\eta_X\mathbf H^X}$ over the background of run $X$, on the blocks of the determining set of $Y$. There is a constant $C_u$ of $d$, $L$, and $G$ such that for $\varrho$ and $\varepsilon_1$ small, at every site $x$ of run~A in a layer $j$ of $Y$,
\begin{equation}\label{eq:gauge-transformations-runs}
    \abs{u^B(x) - u^A(x)} \leq C_u\Bigl(d_{\mathbf H} + \frac{\varepsilon}{a_j}\Bigr)\varrho\;.
\end{equation}
\end{lemma}

\begin{proof}
By \cite[(106)]{BalabanAVG}, for a site $x$ in a block $B^k(y)$ with the sites $x_k = y, x_{k-1}, \dots, x_0 = x$ of the blocks which contain it,
\[
    u(x) = \prod_{j'=k-1}^{0}R\bigl(U_0(\Gamma^{(j'+1)}_{x_{j'+1},x})\bigr)^{-1}\Bigl[\bigl(\overline{\bar R^{j'}_{0,x_{j'+1}}\bar{\bar U}^{j'}_1}\bigr)^{-1}\bigl(\bar R^{j'}_{0,x_{j'+1}}\bar{\bar U}^{j'}_1\bigr)(\Gamma_{x_{j'+1},x_{j'}})\Bigr]\;,
\]
an ordered product over the scales of factors, the factor of the scale $j'$ being the inverse of the block mean \cite[(105)]{BalabanAVG} of the transports of the $j'$-fold average $\bar{\bar U}^{j'}_1$ of the fluctuation \cite[(90), (91)]{BalabanAVG} times its transport along the contour from $x_{j'+1}$ to $x_{j'}$, conjugated by a transport of the background. The formula is printed for one domain, and its use block by block on a determining set with several layers is Ba{\l}aban's \cite[p.~373]{BalabanLFR2}: at a point of a layer $j$ the product runs over the scales $j' < j$.

\emph{The two products.} A site $x$ of run~A is a block point of the first level of run~B, so that run~B's product at $x$ has one factor more, that of its scale $0$, which we call $f^B_0(x)$. By \cite[(91)]{BalabanAVG} the $(j' + 1)$-fold average of run~B is the $j'$-fold average, about $MU^B = V_0$, of its one-step average, and by Proposition~\ref{prop:two-lattices}(j) the one-step average of $e^{ih\mathbf H^B}U^B$ is $e^{i\varepsilon Y}V_0$ with $Y = Q_1\mathbf H^B + c_1$, $c_1 = \varepsilon^{-1}C^B_1(\varepsilon\mathbf H^B)$. Hence the factors of run~B of the scales $j' + 1 \geq 1$ are the factors of run~A of the scales $j'$ evaluated at the field $Y$ in the place of $\mathbf H^A$, up to their prefactors, and
\[
    u^B(x) = u^A[Y](x)\,f^B_0(x)\,(1 + r(x))\;,
\]
where $r$ comes from the prefactors and is bounded below.

\emph{The factor of the finest scale.} By \cite[(82), (87)]{BalabanAVG} at one step, $f^B_0(x)$ is the inverse of a mean of products of at most $dL$ bond variables $e^{ih\mathbf H^B}$ of one block of the side $\varepsilon$, so that $\abs{f^B_0(x) - 1} \leq c\,dLh\sup\abs{\mathbf H^B} \leq c\,d(\varepsilon/a_j)\varrho$ at a point of the layer $j$, because $\abs{\mathbf H^B} \leq \varrho/a_j$ there.

\emph{The dependence on the field.} In the units of the block of the scale $j$ which contains $x$ the field satisfies $\abs A \leq \varrho$, and by \cite[Proposition~4 and (131), pp.~38--39]{BalabanAVG} the angles of the $j'$-fold average are analytic in $A$ for $\abs A < c_4$ and at most $2\alpha_1L^{j'}\eta$ for $\abs A < \alpha_1$, with $\eta = \varepsilon/a_j$; Cauchy's estimate on $\abs A < \frac{1}{2}\alpha_1$ gives them the Lipschitz constant $4L^{j'}\eta$. The factor of the scale $j'$ is the inverse of a mean of at most $L^d$ such angles along contours of at most $dL$ bonds, times the angles along the contour from $x_{j'+1}$ to $x_{j'}$, of at most $dL$ bonds, so that it is Lipschitz continuous in the field with a constant $c\,dL^{j'+1-j}$. The sum over $j' < j$ is at most $c\,dL/(L - 1)$, and since $Y - \mathbf H^A = Q_1(\mathbf H^B - \mathcal L_\iota\mathbf H^A) + c_1$, by $Q_1\mathcal L_\iota = 1$, with $\abs{c_1} \leq C_2\varepsilon\abs{\mathbf H^B}^2$ \cite[(44)]{BalabanVAR},
\[
    \abs{u^A[Y](x) - u^A[\mathbf H^A](x)} \leq c'\bigl(d_{\mathbf H}\varrho + C_2(\varepsilon/a_j)\varrho^2\bigr)\;.
\]

\emph{The prefactors.} The prefactors of \cite[(106)]{BalabanAVG} are transports of the background along the contours: of $U^B$ along the contours of run~B, and of $V_0 = MU^B$ along those of run~A, which are the same lines. The variable $V_0(b)$ of a bond $b$ of run~A is the one-step average of $U^B$ over paths across $b$, each of which differs from the straight one by the flux through an area of at most $d\varepsilon^2$, so that $\abs{\log(V_0(b)^{-1}U^B(\Gamma_b))} \leq c\,d\varepsilon^2\sup\abs{F(U^B)} \leq c\,d\,C\varepsilon_1(\varepsilon/a_j)^2$ by \eqref{eq:minimiser-regularity}. A contour of the scale $j' + 1$ has at most $dL^{j'+1}$ bonds of run~A, so that the two prefactors of that scale differ by at most $c\,d^2C\varepsilon_1(\varepsilon/a_j)L^{j'+1-j}$, and this difference enters multiplied by the deviation of the factor of the scale $j'$ from $1$, which is at most a constant times $L^{j'+1-j}\varrho$. Summed over $j' < j$ this gives $\abs r \leq c''\varepsilon_1(\varepsilon/a_j)\varrho$. The three bounds together are \eqref{eq:gauge-transformations-runs}.
\end{proof}
